\sectionTitle{O Título da seção vem aqui}
Enquanto o texto da seção pode vir aqui, seguido por caixas de destaque.

\begin{mainbox}[center title]{Check out our GitHub Repository \faGithub}
    \begin{center}
        \begin{minipage}[c]{0.2\linewidth}
            \centering
            \qrcode[height=\linewidth]{https://github.com/timotheosf}
        \end{minipage}
    \end{center}
\end{mainbox}

\begin{highlightbox}{A caixa \texttt{highlightbox}}
    Dado um ensemble canônico, definimos sua função de partição $\mathcal{Z}_C$ como a soma dos pesos de Boltzmann sobre \textit{todos} os estados de equilíbrio possíveis ao sistema:
    \begin{equation}
        \mathcal{Z}_C(T,V,N) = \sum_{j}e^{-E(j)\beta}.
    \end{equation}
    Esta função está ligada à energia livre de Helmholtz pela relação:
    \begin{equation}
        F(T,V,N)=-k_BT\ln\mathcal{Z}_C.
    \end{equation}
\end{highlightbox}

Podemos escrever as nossas abobrinhas como bem quisermos.

\blablabla[2]

\begin{mainbox}{A caixa \texttt{mainbox}}
    Dado um ensemble canônico, definimos sua função de partição $\mathcal{Z}_C$ como a soma dos pesos de Boltzmann sobre \textit{todos} os estados de equilíbrio possíveis ao sistema:
    \begin{equation}
        \mathcal{Z}_C(T,V,N) = \sum_{j}e^{-E(j)\beta}.
    \end{equation}
    Esta função está ligada à energia livre de Helmholtz pela relação:
    \begin{equation}
        F(T,V,N)=-k_BT\ln\mathcal{Z}_C.
    \end{equation}
\end{mainbox}

\sectionTitle{Fazer outra seção não quebra o multicol}


E a quebra de coluna não atrapalha as caixas de destaque.


\begin{graybox}{A caixa \texttt{graybox}}
    Dado um ensemble canônico, definimos sua função de partição $\mathcal{Z}_C$ como a soma dos pesos de Boltzmann sobre \textit{todos} os estados de equilíbrio possíveis ao sistema:
    \begin{equation}
        \mathcal{Z}_C(T,V,N) = \sum_{j}e^{-E(j)\beta}.
    \end{equation}
    Esta função está ligada à energia livre de Helmholtz pela relação:
    \begin{equation}
        F(T,V,N)=-k_BT\ln\mathcal{Z}_C.
    \end{equation}
\end{graybox}

\begin{darkbox}{A caixa \texttt{darkbox}}
    Dado um ensemble canônico, definimos sua função de partição $\mathcal{Z}_C$ como a soma dos pesos de Boltzmann sobre \textit{todos} os estados de equilíbrio possíveis ao sistema:
    \begin{equation}
        \mathcal{Z}_C(T,V,N) = \sum_{j}e^{-E(j)\beta}.
    \end{equation}
    Esta função está ligada à energia livre de Helmholtz pela relação:
    \begin{equation}
        F(T,V,N)=-k_BT\ln\mathcal{Z}_C.
    \end{equation}
\end{darkbox}


Podemos dar \Highlight{destaque a algum texto} com o comando \verb|\Highlight|.

E, no final, temos os agradecimentos e referências impressos automaticamente, com suporte às citações \texttt{natbib}.
